% Meta-Informationen -------------------------------------------------------
%		Informationen über das Dokument, wie z.B. Titel, Autor, Matrikelnr. etc
%		werden in der Datei _Meta.tex definiert und können danach global
%		verwendet werden.
% --------------------------------------------------------------------------
% Informationen ------------------------------------------------------------
% 	Definition von globalen Parametern, die im gesamten Dokument verwendet
% 	werden können (z.B auf dem Deckblatt etc.).
% --------------------------------------------------------------------------
\newcommand{\titel}{Ermittlung des stationären Zustands beim Festplattenbenchmarking}
\newcommand{\art}{Bachelorarbeit} %Bachelorarbeit
\newcommand{\ort}{Leipzig}
\newcommand{\hochschule}{Universität Leipzig}
\newcommand{\fachgebiet}{Abteilung Datenbanken}
\newcommand{\fakultaet}{Fakultät für Mathematik und Informatik}
\newcommand{\institut}{Institut für Informatik}
\newcommand{\autor}{Nick Meier}
\newcommand{\matrikelnr}{3712562}
\newcommand{\erstbetreuer}{Prof. Dr. Erhard Rahm}
\newcommand{\zweitbetreuer}{Dr. David Georg Reichelt}
\newcommand{\jahr}{2025}
\newcommand{\invnr}{1337}
\newcommand{\eingereicht}{xx.xx.xxxx}

% Eigene Befehle
\newcommand{\todo}[1]{\textbf{\textsc{\textcolor{red}{(TODO: #1)}}}}

% Autorennamen in small caps
\newcommand{\AutorZ}[1]{\textsc{#1}}
\newcommand{\Autor}[1]{\AutorZ{\citeauthor{#1}}}

% Befehle zur semantischen Auszeichnung von Text
\newcommand{\NeuerBegriff}[1]{\textbf{#1}}
\newcommand{\Fachbegriff}[1]{\textit{#1}}
\newcommand{\Prozess}[1]{\textit{#1}}
\newcommand{\Webservice}[1]{\textit{#1}}
\newcommand{\Eingabe}[1]{\texttt{#1}}
\newcommand{\Code}[1]{\texttt{#1}}
\newcommand{\Datei}[1]{\texttt{#1}}
\newcommand{\Datentyp}[1]{\textsf{#1}}
\newcommand{\XMLElement}[1]{\textsf{#1}}

% Abkürzungen
\newcommand{\vgl}{Vgl.\ }
\newcommand{\ua}{\mbox{u.\,a.\ }}
\newcommand{\zB}{\mbox{z.\,B.\ }}
\newcommand{\bs}{$\backslash$}

% Einfache Anführungszeichen in texttt
\newcommand{\sq}{\textquotesingle}



% Dokumentenkopf -----------------------------------------------------------
% 	Diese Vorlage basiert auf "scrreprt" aus dem koma-script.
%		Die Option draft sollte beim fertigen Dokument ausgeschaltet werden.
% --------------------------------------------------------------------------
\documentclass[
	11pt,					% Schriftgröße
	DIV=10,
	ngerman,				% für Umlaute, Silbentrennung etc.
	a4paper,				% Papierformat
	oneside,				% einseitiges Dokument
	titlepage,				% es wird eine Titelseite verwendet
	parskip=half,			% Abstand zwischen Absätzen (halbe Zeile)
	headings=normal, % Größe der Überschriften verkleinern
	numbers=withendperiod, % Fügt in den Überschriften nach den Zahlen einen Punkt ein
	listof=totoc,				% Verzeichnisse im Inhaltsverzeichnis aufführen
	bibliography=totoc,				% Literaturverzeichnis im Inhaltsverzeichnis aufführen
	index=totoc,				% Index im Inhaltsverzeichnis aufführen
	captions=tableheading,		% Beschriftung von Tabellen oberhalb ausgeben
	final					% Status des Dokuments (final/draft)
]{scrreprt}

\renewcommand*\chapterheadstartvskip{\vspace*{-1.0cm}}

% Bentigte Packages -------------------------------------------------------
%		Weitere Packages, die benötigt werden, sind in die Datei Packages.tex
%		"ausgelagert", um die Vorlage möglichst übersichtlich zu halten.
% --------------------------------------------------------------------------
% Anpassung des Seitenlayouts ----------------------------------------------
% 	siehe Seitenstil.tex
% --------------------------------------------------------------------------
\usepackage[
	automark,			% Kapitelangaben in Kopfzeile automatisch erstellen
	headsepline,	% Trennlinie unter Kopfzeile
	ilines				% Trennlinie linksbündig ausrichten
]{scrlayer-scrpage}
\usepackage{scrhack} % Disable some warnings

\usepackage{pseudocode}
\usepackage{nicefrac}

% Für eine schöne Anordnung von Bildern
\usepackage{subfigure}

\usepackage{dsfont}
%\usepackage{color}
%
%% Define user colors using the RGB model
%\definecolor{yellow}{rgb}{0.0,1.0,0.0}
%\definecolor{rot}{rgb}{1.0,0.0,0.0}

% Anpassung an Landessprache -----------------------------------------------
% 	Verwendet globale Option german siehe \documentclass
% --------------------------------------------------------------------------
\usepackage[ngerman]{babel}

% Umlaute ------------------------------------------------------------------
% 		Umlaute/Sonderzeichen wie äöüß direkt im Quelltext verwenden (CodePage).
%		Erlaubt automatische Trennung von Worten mit Umlauten.
% --------------------------------------------------------------------------
\usepackage[utf8]{inputenc}
\usepackage[T1]{fontenc}
%\usepackage{ae} % "schöneres" ä
\usepackage{textcomp} % Euro-Zeichen etc.
\usepackage{lmodern} % schööön

% Grafiken -----------------------------------------------------------------
% 		Einbinden von Grafiken [draft oder final]
% 		Option [draft] bindet Bilder nicht ein - auch globale Option
% --------------------------------------------------------------------------
\usepackage[dvips,final]{graphicx}
\usepackage{wrapfig}
\graphicspath{{Bilder/}} % Dort liegen die Bilder des Dokuments

% Befehle aus AMSTeX für mathematische Symbole z.B. \boldsymbol \mathbb ----
\usepackage{amsmath,amsfonts,amsthm}

% Für Index-Ausgabe; \printindex -------------------------------------------
\usepackage{makeidx}

% Einfache Definition der Zeilenabstände und Seitenränder etc. -------------
\usepackage{setspace}
\usepackage{geometry}

% für gedrehte Tabellen
\usepackage{rotating} 

% Symbolverzeichnis --------------------------------------------------------
% 	Symbolverzeichnisse bequem erstellen, beruht auf MakeIndex.
% 		makeindex.exe %Name%.nlo -s nomencl.ist -o %Name%.nls
% 	erzeugt dann das Verzeichnis. Dieser Befehl kann z.B. im TeXnicCenter
%		als Postprozessor eingetragen werden, damit er nicht ständig manuell
%		ausgeführt werden muss.
%		Die Definitionen sind ausgegliedert in die Datei Abkuerzungen.tex.
% --------------------------------------------------------------------------
\usepackage[intoc]{nomencl}
  \let\abbrev\nomenclature
  \renewcommand{\nomname}{Abkürzungsverzeichnis}
  \setlength{\nomlabelwidth}{.25\hsize}
  \renewcommand{\nomlabel}[1]{#1 \dotfill}
  \setlength{\nomitemsep}{-\parsep}

% Zum Umfließen von Bildern -------------------------------------------------
\usepackage{floatflt}

% Zum Einbinden von Programmcode --------------------------------------------
\usepackage{listings}
\usepackage{xcolor} 
\definecolor{hellgelb}{rgb}{1,1,0.9}
\definecolor{colKeys}{rgb}{0,0,1}
\definecolor{colIdentifier}{rgb}{0,0,0}
\definecolor{colComments}{rgb}{1,0,0}
\definecolor{colString}{rgb}{0,0.5,0}
\lstset{%
    float=hbp,%
    basicstyle=\texttt\small, %
    identifierstyle=\color{colIdentifier}, %
    keywordstyle=\color{colKeys}, %
    stringstyle=\color{colString}, %
    commentstyle=\color{colComments}, %
    columns=flexible, %
    tabsize=2, %
    frame=single, %
    extendedchars=true, %
    showspaces=false, %
    showstringspaces=false, %
    numbers=left, %
    numberstyle=\tiny, %
    breaklines=true, %
    backgroundcolor=\color{hellgelb}, %
    breakautoindent=true, %
%    captionpos=b%
}

% Lange URLs umbrechen etc. -------------------------------------------------
\usepackage{url}


%% Wichtig für korrekte Zitierweise ------------------------------------------

\usepackage[autocite=inline, sorting=none]{biblatex}
\bibliography{quellen} % Name der .bib-Datei

\usepackage{csquotes} % Empfohlen, um Zitierten Text richtig darzustellen

% ermöglicht Zeilenumbrüche in Captions
\usepackage{caption}


% PDF-Optionen --------------------------------------------------------------
\usepackage[
bookmarks,
bookmarksopen=true,
pdftitle={\titel},
pdfauthor={\autor},
pdfcreator={\autor},
pdfsubject={\titel},
pdfkeywords={\titel},
colorlinks=true,
%linkcolor=red, % einfache interne Verknüpfungen
%anchorcolor=black,% Ankertext
%citecolor=blue, % Verweise auf Literaturverzeichniseinträge im Text
%filecolor=magenta, % Verknüpfungen, die lokale Dateien öffnen
%menucolor=red, % Acrobat-Menüpunkte
%urlcolor=cyan, 
% für die Druckversion können die Farben ausgeschaltet werden:
linkcolor=black, % einfache interne Verknüpfungen
anchorcolor=black,% Ankertext
citecolor=black, % Verweise auf Literaturverzeichniseinträge im Text5
filecolor=black, % Verknüpfungen, die lokale Dateien öffnen
menucolor=black, % Acrobat-Menüpunkte
urlcolor=black, 
%backref,
%pagebackref,
plainpages=false,% zur korrekten Erstellung der Bookmarks
pdfpagelabels,% zur korrekten Erstellung der Bookmarks
hypertexnames=false,% zur korrekten Erstellung der Bookmarks
linktocpage % Seitenzahlen anstatt Text im Inhaltsverzeichnis verlinken
]{hyperref}

% Zum fortlaufenden Durchnummerieren der Fußnoten ---------------------------
\usepackage{chngcntr}


% für lange Tabellen
\usepackage{longtable}
\usepackage{array}
\usepackage{tabularx}
\usepackage{ragged2e}
\usepackage{lscape}

\usepackage{supertabular}


% Tables mit punkten

\usepackage{multirow}
\usepackage{booktabs}

% Spaltendefinition rechtsbündig mit definierter Breite ---------------------
\newcolumntype{w}[1]{>{\raggedleft\hspace{0pt}}p{#1}}

% Formatierung von Listen ändern
\usepackage{paralist}
% Standardeinstellungen:
% \setdefaultleftmargin{2.5em}{2.2em}{1.87em}{1.7em}{1em}{1em}

\usepackage{tablefootnote}
% für Ausblenden der Seitenzahl
\usepackage{lipsum}

% Erstellung eines Index und Abkürzungsverzeichnisses aktivieren -----------
\makeindex
% makeindex Masterarbeit.nlo -s nomencl.ist -o Masterarbeit.nls
\makenomenclature


% Kopf- und Fußzeilen, Seitenränder etc. -----------------------------------
% Zeilenabstand ------------------------------------------------------------
\onehalfspacing 
% \setstretch{1,5}

% Seitenränder -------------------------------------------------------------
\geometry{paper=a4paper,left=25mm,right=20mm,top=20mm, bottom=25mm}
% Notfall maße :)
%\geometry{paper=a4paper,left=35mm,right=25mm,top=25mm, bottom=25mm}



% Kopf- und Fußzeilen ------------------------------------------------------
\pagestyle{scrheadings}

% Kopf- und Fußzeile auch auf Kapitelanfangsseiten -------------------------
\renewcommand*{\chapterpagestyle}{scrheadings}

% Schriftform der Kopfzeile ------------------------------------------------
\renewcommand{\headfont}{\normalfont}

% Kopfzeile ----------------------------------------------------------------
\ihead{\textit{\headmark}}
\chead{}
%\ohead{\includegraphics[scale=1]{Bilder/logoKlein.JPG}}
\ohead{}
\setlength{\headheight}{8mm} % Höhe der Kopfzeile
\setheadwidth[0pt]{textwithmarginpar} % Kopfzeile über den Text hinaus verbreitern

% Fußzeile -----------------------------------------------------------------
% \ifoot{\copyright\ \autor \\ \invnr}
% \ifoot{\copyright\ \autor \\ \matrikelnr}
\ifoot{\autor \\ \matrikelnr}
\cfoot{}
\ofoot{\pagemark}
\setlength{\footskip}{12mm}
\setfootwidth[0pt]{text}

% Überschriften ------------------------------------------------------------
\renewcommand*\chapterheadstartvskip{\vspace*{-0.5cm}} % Platz vor einer Überschrift eines neuen Kapitels


% erzeugt ein wenig mehr Platz hinter einem Punkt --------------------------
\frenchspacing

% Schusterjungen und Hurenkinder vermeiden
\clubpenalty = 10000
\widowpenalty = 10000 
\displaywidowpenalty = 10000


% Quellcode-Ausgabe formatieren --------------------------------------------
%\lstset{numbers=left, numberstyle=\tiny, numbersep=5pt, breaklines=true}
%\lstset{emph={square}, emphstyle=\color{red}, emph={[2]root,base}, emphstyle={[2]\color{blue}}}

\definecolor{gray}{rgb}{0.9,0.9,0.9}

\lstset{%
		basicstyle=\small\ttfamily,language={[LaTeX]TeX},
		numbersep=5mm, 
		numbers=left,
		numberstyle=\tiny,
		breaklines=true,
		framexleftmargin=8mm, 
		xleftmargin=8mm,
		backgroundcolor=\color{gray},
		captionpos=b
}%

% Fußnoten fortlaufend durchnummerieren ------------------------------------
\counterwithout{footnote}{chapter}

% Definitionen

\newtheorem{definition}{Definition}


% Eigene Definitionen für Silbentrennung
\hyphenation{Trenn-bar-es}

% Das eigentliche Dokument -------------------------------------------------
%		Der eigentliche Inhalt des Dokuments beginnt hier. Die einzelnen Seiten
%		und Kapitel werden in eigene Dateien ausgelagert und hier nur inkludiert.
% --------------------------------------------------------------------------

\begin{document}
% auch subsubsection nummerieren
\setcounter{secnumdepth}{3}
\setcounter{tocdepth}{3}

% keine Kopf-/Fußzeilen bei Deckblatt und Abstract
\ofoot{}
% Deckblatt
\thispagestyle{plain}
\begin{titlepage}

\begin{center}
\includegraphics[height=7cm]{Bilder/Uni-L.png}\\[2.5ex]

\institut\\
\fakultaet\\
\fachgebiet\\[6ex]

\textbf{\large\titel}\\[1.5ex]
\art\\[6ex]

\normalsize
vorgelegt von:\\
\autor\\[1.5ex]
Matrikelnummer:\\
\matrikelnr\\[1.5ex]
Betreuer:\\
\erstbetreuer\\
\zweitbetreuer\\[1.0ex]
\end{center}

%\begin{tabbing}
%\hspace{3.5cm}\= \kill
%   vorgelegt von: \> \autor\\[1.2ex]
%   Matrikelnummer: \> \matrikelnr\\[1.2ex]
%    \> \\
%   Betreuer: \> \erstbetreuer\\[1.2ex]
%    \> \zweitbetreuer
%\end{tabbing}

\begin{center}
\copyright\ \jahr\\[1.0ex]
\end{center}

\singlespacing
\small
\noindent Dieses Werk einschließlich seiner Teile ist \textbf{urheberrechtlich geschützt}. Jede Verwertung außerhalb der engen Grenzen des Urheberrechtgesetzes ist ohne Zustimmung des Autors unzulässig und strafbar. Das gilt insbesondere für Vervielfältigungen, Übersetzungen, Mikroverfilmungen sowie die Einspeicherung und Verarbeitung in elektronischen Systemen.

\end{titlepage}


\section*{Zusammenfassung}
\label{sec:Zusammenfassung}

Mit Benchmarks kann man die Performanz von Hardware oder Software, wie etwa die Rechenleistung und die I/O Geschwindigkeit messen. 
Diese Messungen können stark oder schwach schwanken. Dies hängt unter anderem davon ab, ob die Tests gerade erst begonnen haben oder über einen längeren Zeitraum ausgeführt wurden.
Ein stationärer Zustand ist ein Muster, das sich mit der Zeit wiederholt und bei dem das Verhalten in einem Zeitraum von der gleichen Natur ist wie in jedem anderen Zeitraum.
Um zu ermitteln, wann der stationäre Zustand eingetreten ist, sind stabile Messwerte mit minimalen Schwankungen erforderlich.
Das Problem ist besonders bei der Netzwerkübertragung relevant, da zusätzlich die Netzwerkübertragung die Schwankung beeinflusst.
Diese Messungen der I/O-Geschwindigkeit werden mit dem Benchmark-Tool Fio durchgeführt, und es wird ein Programm entwickelt, das diesen Zustand analysiert. 

In dieser Arbeit wird die statistische Analyse durchgeführt, um zu ermitteln, ab wann sich Messzeiträume nicht mehr statistisch signifikant unterscheiden. Die Tests, die dafür verwendet werden, sind der t-Test, Tukey-HSD-Test und der Mann-Whitney-Test. 
Basierend auf der Analyse kann eine Aussage darüber getroffen werden, wie oft die fio-Messungen wiederholt werden müssen, um den stationären Zustand beim Festplattenbenchmarking zu erreichen.
Das Programm wird in Java implementiert.
% \section*{Danksagung}
\label{sec:Danksagung}
Danksagung
\newpage
\ofoot{\pagemark}

% Seitennummerierung -------------------------------------------------------
%		Vor dem Hauptteil werden die Seiten in großen römischen Ziffern 
%		nummeriert...
% --------------------------------------------------------------------------
\pagenumbering{Roman}

\tableofcontents			% Inhaltsverzeichnis

% Abkürzungsverzeichnis ----------------------------------------------------
%\input{Inhalt/Glossar}
%\printnomenclature
%\label{sec:Glossar}

\listoffigures					% Abbildungsverzeichnis
\listoftables					% Tabellenverzeichnis

%\renewcommand{\lstlistlistingname}{Verzeichnis der Listings}
%\lstlistoflistings	

% ...danach in normalen arabischen Ziffern ---------------------------------
\clearpage
\pagenumbering{arabic}


% Inhalt -------------------------------------------------------------------
%		Hier können jetzt die einzelnen Kapitel inkludiert werden. Sie müssen
%		in den entsprechenden .TEX-Dateien vorliegen. Die Dateinamen können
% 		natürlich angepasst werden.
% --------------------------------------------------------------------------
\chapter{Einleitung}
\label{cha:Einleitung}
1 Einf¨uhrung
Die Performanz eines Rechners ist die Leistungsf¨ahigkeit und Effizienz bei der
Ausf¨uhrung von Aufgaben [5]. Um herauszufinden welche Systeme die beste
Performanz liefert, sind Messungen notwendig, um sie vergleichen zu k¨onnen.
Solche Messungen und Vergleiche werden mit Benchmarks durchgef¨uhrt [14].
Mit ihnen wird die Rechenleistung oder auch Lese-/Schreibgeschwindigkeit der
Hardware oder auch Software mit standardisierten Tests gemessen [4]. Ein
Beispiel f¨ur getestete Software, sind virtuelle Maschinen wie die Java Virtual
Machine (JVM) [6] oder virtuelle Maschinen die reale Rechnern simulieren [9].
Durch Benchmarks wird die Performanz von virtuellen CPUs oder des Hyper-
visor analysiert. Auch f¨ur Cloud Computing Plattformen wird Benchmarking
betrieben, um auch Performanz zwischen verschieden Cloud-Plattformen ver-
gleichen zu k¨onnen. Diese Benchmarks verwenden z.B die Turnaround-Zeit
als Metrik, also die Zeit, die ein Prozess ben¨otigt, um vollst¨andig durchzu-
laufen [10]. Somit gibt es ein weites Spektrum an Auswahl von Systemen, wo
Benchmarking betrieben wird. Ein Beispielprogramm f¨ur so ein Benchmark Tool
ist das Fio1. Fio ist ein Tool, mit dem man die Bandbreitengeschwindigkeit vom
Lesen/Schreiben testet und sie als eine Textdatei (von Fio Logdatei genannt)
ausgeben kann. Diese Tests werden in Fio Jobs genannt.
Die I/O-Geschwindigkeit kann st¨arker oder schw¨acher vom tats¨achlich er-
reichbaren Maximalwert abweichen. Das Ziel des geplanten Tools ist es, den
station¨aren Zustand zu ermitteln. Ein Zustand, wo die Rechenleistung oder
Lese-/Schreibge-schwindigkeit sich nicht mehr beim Messen stark ¨andert und
sich einem konstanten Wert ann¨ahert, aber durch den Nichtdeterminismus nie
erreicht wird [3]. Um den station¨aren Zustand zu ermitteln, werden verschiedene
Methodiken aus der Varianzanalyse verwendet. Der Tukey-HSD-Test, t-Test
und Mann-Whitney-Test sind statistische Tests die im Analysetool genutzt wer-
den. Zusammen mit den Logs aus den Jobs des Fio’s wird die Standabweichung
sowie die Varianz berechnet. Dabei werden mit den Tests die statistische Sig-
nifikanz als auch den Zusammenhang zwischen den Jobs analysiert.
Im Expos´e werden zuerst die Grundlagen behandelt. Es werden das Fio-Tool
sowie die Definitionen des station¨aren Zustands und die verwendeten statistis-
chen Tests beschrieben. Im zweiten Teil wird das Konzept beschrieben, wie
das entwickelte Programm mit dem Fio-Tool arbeiten wird und es werden erste
Experimente durchgef¨uhrt. Im vorletzten und letzten Kapitel werden auf ver-
wandte Arbeiten eingegangen und es wird der Zeitplan beschrieben, wie viele
Wochen jeweils die Kapitel f¨ur die Bachelor Arbeit genutzt werden.
\section{Motivation}
\section{Ansatz}
\section{Aufbau der Arbeit}
\chapter{Grundlagen}
\label{cha:Grundlagen}
Benchmarking spielt eine wichtige Rolle zum Testen von Hardware und Software.
Um die Performanz von Hardware [Q] oder auch Software [Q] zu untersuchen Tests verwendet
um Performanz wie I/O-Geschwindigkeit oder CPU-Geschwindigkeit zu messen [Q].


\section{Benchmarking}

Benchmarks sind Tests, die die Performanz von Hardware oder Software misst. Diese Messungen werden evaluiert
um verschiedene Hardware zu vergleichen [https://www.sciencedirect.com/topics/computer-science/performance-benchmark].
Das verwendete Benchmarking-Tool fio misst die Performanz der I/O-Geschwindigkeit einer Festplatte oder auch eines Dateisystems. [fio]
Probleme die bei den Messungen begrenzen sind nicht-deterministische Faktoren wie buffer caches, disk caches oder auch kernel daemons [A Nine Year Study].

\subsection{Fio-Tool}
Das Benchmarking-Tool arbeitet in der Konsole, z.B in Cygwin, und es gibt keine UI.
Um das Programm auszuführen werden zusätzlich Parameter verwendet.
Um einen Festplattenbenchmark-Test selber durchzuführen, sieht der Command wie folgt aus:

$./fio --rw=write --runtime=2s --write_bw_log=mytest --name=test --size=1024m$

\begin{center}
  \begin{table}
    \begin{tabularx}{\textwidth}{|X|X|}
      \hline
        Parameter& Defintion \\ 
      \hline
      fio & Das fio Executable  \\ 
      \hline
      rw & Auswahl von read oder write des Tests  \\ 
      \hline
      runtime & Maximale Dauer eines Tests  \\ 
      \hline
      write bw log &  Name des Logs die ausgeben werden soll   \\ 
      \hline
      name &  Name des Tests   \\ 
      \hline
      size & Größe der Datei für den read/write Test    \\ 
      \hline
    \end{tabularx}
    \caption{Beispiel Command zur Ausführung des Fio-Tools}
    \label{tab:1d_1_sta}
  \end{table}
\end{center}

Die Commands können auch als eine .fio geschrieben werden. In dieser Arbeit so eine Ausführung als einen \textit{run/job} bezeichnet.
Der Job oben testet das Random Read mit einer Laufzeit (Runtime) von 2 Sekunden und liest eine Datei mit einer Größe von 128 Mebibyte.

\subsection{Warmup und stationärer Zustand}

Bevor der stationäre Zustand erfolgt, ist ein Warmup erforderlich.
Dieser Warmup wird auch als transiente Zustand bezeichnet.

Bei Hardware oder Software ist anfangs wiederholtes testen notwendig, um
sie analysieren zu können. Erst nach diesen sogenannten Warm-ups wird der
station¨are Zustand der Performanz (engl. steady state of peak performance)
erreicht [3].

Wenn man das Lesen/Schreiben auf demselben System mit derselben Datei
testet, ist die Bandbreitengeschwindigkeit im Durchschnitt nie die gleiche. Sie
wird immer abweichen. Ursachen daf¨ur könnten schon verschiedene CPU-Temperaturen sein, 
aber auch CPU-Scaling oder parallele Prozessierung [7]. Das Fio
selbst arbeitet nicht nur mit einem Thread, sondern es arbeitet mit Multi-
threads, was nicht-deterministische Zustände hervorrufen kann. Weitere nicht-
deterministische Faktoren sind die Caches die beim Lesen oder Schreiben ver-
wendet werden, je nach dem ob sie leer sind oder schon beschrieben wurden.
Auch nicht essentielle Prozesse (oder auch essentielle Prozesse wie Daemons3),
die immer Hintergrund laufen, können die Geschwindigkeit beeinflussen [15].

Der stationärer Zustand ist ein Muster, das sich mit der Zeit wiederholt und bei dem das Verhalten in einem Zeitraum von der 
gleichen Natur ist wie in jedem anderen Zeitraum. Um zu ermitteln,
wann der stationäre Zustand eingetreten ist, sind stabile Messwerte erforderlich.
Da der Nichtdeterminismus Schwankungen in den Messungen verursacht, müssen dies in Betracht gezogen werden.
Diese Schwankung die bei den Messungen entstehen werden mit einer selbst gewählt Fehlerrate innerhalb des Analysetools ausgewählt.

\textit{Der Standardwert wird $\alpha$ = 0.05 sein.}. Die Fehlerrate dient für die Ermittlung des stationäre Zustandes.
Da sich die Geschwindigkeit nie konstant einem Wert nähert, sondern stets
abweicht, soll das Analysetool in Zukunft mit statistischen Tests mit der zusätzlichen Fehlerrate den stationären
Zustand ermitteln und dabei den Nichtdeterminismus berücksichtigen.

 Ein Beispiel, wo solche Warmups erforderlich sind, ist die JIT
(Just-In-Time) Kompilierung in Java. Der Bytecode wird während der Laufzeit
des Programms kompiliert. Damit verbessert man die Performanz von Java-
Anwendungen [8]. Ein stationärer Zustand im Festplattenbenchmarking ist er-
reicht, wenn die I/O-Geschwindigkeit wiederholte Muster aufweist, sodass die
Geschwindigkeiten zu verschiedenen Zeitpunkten nach diesem Zustand gleich
sind.2 Vor dem stationären Zustand erfolgt der Warm-Up, wo die Caches des
Rechners befüllt werden [15]. Dieser Zeitraum wird auch als transienten Zustand
bezeichnet [11]. Auch wenn der Warm-up erfolgt ist, ist der Nichtdeterminismus
dabei ein Hindernis.

Die Abweichungen können im niedrigen Prozentbereich liegen. Es wird jedoch nicht
ausreichen, ein paar Tests durchzuführen und anschließend die Logs auszuwerten. 
Denn nach dem Warm-up wird ein größeres Problem darin bestehen,
welche Logs die Informationen beinhalten, wo auch der Warm-Up stattfand und
der transiente Zustand abgeschlossen wurde. Die Auswahl der Logs und deren
Auswertung können signifikante Unterschiede aufweisen, die die Evaluierung
verändern würden [1].



\section{Statistische Analyse/Tests}
Um zu entscheiden, ob eine Performanzänderung vorliegt, müssen durch die Messung
ermittelte Stichproben von Zufallsvariablen verglichen werden. Hierfür werden verschie-
dene statistische Tests eingesetzt [GBE07]. 
Die Grundgesamtheit werden alle Logs sein und daraus werden paarweise die Stichproben entnommen die für die Tests verwendet.
Diese statistische Tests wird ausgesagt, ob eine statistische Signifikanz zwischen den Stichproben vorhanden ist.
Wenn die statistisch Signifikanz erfüllt ist, wird die Nullhypothese abgelehnt.
Die Nullhypothese ist hier, zwischen den Logs keine Abhängigkeit besteht.
Somit ist die Alternativhypothese, dass zwischen den ausgewählten Stichproben ein Zusammenhang besteht.
Um die Nullhypothese zu überprüfen wird ein p-wert (Signifikanz Wahrscheinlichkeit) berechnet und mit einem selbst
festgelegtem Signifikanzniveau $\alpha$ verglichen.
Dieser p-Wert sagt aus wie die Wahrscheinlichkeit ist eine Nullhypothese falsch anzunehmen. 
Ein p-Wert von 0.5 sagt aus, dass es eine 5\%-ige Wahrscheinlichkeit gibt das die Nullhypothese inkorrekt abgelehnt wurde, obwohl sie erfüllt ist.
Somit ist die Wahrscheinlichkeit für die Alternativhypothese (1 - p) also 95\% \cite{statistics}.
Wenn $\alpha > p$ mit $\alpha = 0.05$ ist eine statistische Signifikanz vorhanden ist, wird die Nullhypothese abgelehnt. 
Ein weiterer Messwert in den Tests ist der z-Wert. Dieser Wert gibt an viele Standardabweichungen ein ROHwert unter
bzw. über dem Mittelwert liegt.

Im Folgenden werden statistische Tests im
Allgemeinen, der T-Test, der Konfidenzintervallvergleich und weitere statistische Tests
vorgestellt.

\section{Varianzanalyse - Analysis of Variance (ANOVA)}

Die ANOVA dient als statistische Analyse der Varianz und Mittelwert mit mehreren Auswahl von Stichproben.




\begin{table}[h!]
  \centering
  %\resizebox{\textwidth}{!}{
  \begin{tabular}{|c|*{5}{c}|c|}
  \hline
  \textbf{Messwerte} & \multicolumn{5}{c|}{\textbf{Iterationen}} & \textbf{Overall mean} \\
  \cline{2-6}
   & 1 & 2 & $\cdots$ & $j$ & $k$ & \\
  \hline
  1 & $y_{11}$ & $y_{12}$ & $\cdots$ & $y_{1j}$ & $\cdots$ & $y_{1k}$ \\
  2 & $y_{21}$ & $y_{22}$ & $\cdots$ & $y_{2j}$ & $\cdots$ & $y_{2k}$ \\
  $\vdots$ & $\vdots$ & $\vdots$ & $\ddots$ & $\vdots$ & $\ddots$ & $\vdots$ \\
  $i$ & $y_{i1}$ & $y_{i2}$ & $\cdots$ & $y_{ij}$ & $\cdots$ & $y_{ik}$ \\
  $\vdots$ & $\vdots$ & $\vdots$ & $\ddots$ & $\vdots$ & $\ddots$ & $\vdots$ \\
  $n$ & $y_{n1}$ & $y_{n2}$ & $\cdots$ & $y_{nj}$ & $\cdots$ & $y_{nk}$ \\
  \hline
  Column means & $\bar{y}_{\cdot1}$ & $\bar{y}_{\cdot2}$ & $\cdots$ & $\bar{y}_{\cdot j}$ & $\cdots$ & $\bar{y}_{\cdot k}$ \\
  Effects & $\alpha_1$ & $\alpha_2$ & $\cdots$ & $\alpha_j$ & $\cdots$ & $\alpha_k$ \\
  \hline
  \end{tabular}
  %}
  \caption{Tabelle für einen fio-Run mit k-Iterationen}
  \label{tab:measurements}
\end{table}

\section{Konfidenzintervalle}
Es wird angenommen das die Grundgesamtheit normalverteilt sind [QUELLE].
Die Stichprobenverteilung wird in eine Standnormalverteilung überführt, mit:

\begin{center}
  $z = \dfrac{\overline{x} - \mu}{\sigma_{\overline{x}}}$
\end{center}

Die Standardabweichung sagt nicht aus ob der Wert kein "ausreißer" ist. Dafür wird die Konfidenzintervalle verwendet
um akkurater zu sein.
Die Wahrscheinlichkeit dafür, dass der z-Wert
größer ist als das $z_{\alpha/2}$ und kleiner
als das $z_{1-\alpha/2}$ Perzentil ist, beträgt $1 - \alpha$:

\begin{center}
  $P(z_{\alpha/2} \leq z \leq z_{1-\alpha/2}) = 1 - \alpha$
\end{center}

Mit hilfe der obigen Formel lässt sich diese Formel wie folgt umformen:

\begin{center}
  $P(\overline{x} - z_{1-\alpha/2} * \sigma_{\overline{x}} \leq \mu \leq \overline{x} + z_{1-\alpha/2} * \sigma_{\overline{x}}) = 1 - \alpha$
\end{center}

Diese Wahrscheinlichkeitsformel gibt an, wie der $\mu$-Wert sich in einem Konfidenzintervall mit $(1 - \alpha)$ befindet.
Der Bereich lässt sich wie folgt berechnen:

\begin{center}
  Konfidenzintervall $= [\overline{x} - z_{1-\alpha / 2} \cdot \sigma_{\overline{x}}, \overline{x} + z_{1-\alpha / 2} \cdot \sigma_{\overline{x}}]$
\end{center}

\subsection{T-Test}
Mit dem T-Test werden jeweils zwei unabhängige Stichprobe miteinander verglichen, 
unter der Annahme das die Grundgesamtheit normalverteilt ist.

\begin{center}
  $t = \dfrac{\mu_1-\mu_2}{\sqrt{\frac{s_p^2}{n_1} + \frac{s_p^2}{n_2}}}$
\end{center}

\begin{center}
  $s^2_p = \dfrac{(n_1 - 1)s^2_1 + (n_2 - 1)s^2_2}{n_1 + n_2 - 2}$
\end{center}

\begin{center}
  $|t| > t(1 - \dfrac{1}{2}\alpha, n + m - 2)$
\end{center}

Der Begriff T-Test wird genutzt, um verschiedene Tests zu bezeichnen, die
u. a. die Verteilung der Grundgesamtheit von normalverteilten und unabhängigen Stich-
proben miteinander vergleichen [JP15, Seite 41]. Im Kontext dieser Arbeit bezeichnet
der T-Test den unabhängigen T-Test, bei dem die Erwartungswerte zweier voneinander
unabhängiger Stichproben miteinander verglichen werden.
Der T-Test vergleicht die Stärke der Unterschiede zwischen den Stichproben mit der
Abweichung innerhalb der Stichproben. Die Prüfgröße t, die ein Maß dafür ist, wie un-
wahrscheinlich die Nullhypothese H0 ist, d. h. dass die Grundgesamtheiten der beiden
Stichproben den gleichen Erwartungswert aufweisen. Diese wird folgendermaßen berech-
net:
\subsection{Tukey-HSD}
Der Tukey-HSD-Test verwendet die Differenz zwischen zwei Mit-
telwerten, um ihre statistische Signifikanz zu ermitteln.
Der HSD-Wert wird mittels mittlere quadratische Abweichung (MSE) berech-
net und mit der Anzahl der Gruppen die getestet wurden. Die minimale Dif-
ferenz wird als HSD bezeichnet. Wenn aber die Differenz der Mittelwerte aus
den Tests gr¨oßer als der HSD-Wert ist, sind diese Werte statistisch signifikant.
Zus¨atzlich versucht der Tukey-HSD-Test5 die Fehlerrate von Falsch-Positiven-
Annahmen zu korrigieren. Dieser Test setzt aber voraus, dass die Stichproben
normalverteilt sind. Die t-Tests arbeiten ¨ahnlich wie die Tukey HSD Tests. Der
Unterschied besteht darin, dass sie die Fehlerrate nicht anpassen [12]. 

\subsection{Mann-Whitney-
Test}
Der Mann-Whitney-Wilcoxon-Rangsummentest [BBH+16; CGT+14; JHHF09]
[Fie13, Seiten 541ff] kann ebenfalls eingesetzt werden, um zu entscheiden, ob sich zwei
Stichproben sich signifikant unterscheiden in dem die Werte einen Rang gesetzt wird. 
Alle Werte aus den Stichproben bekommen einen Rang aufsteigend vom Kleinsten zum Größten Wert.
Die gesamte Rangsumme lässt mit der Formel berechnen:
\begin{center}
  $T_1 + T_2 = \dfrac{n * (n + 1)}{2}, (n = n_1 + n_2)$
\end{center}  

Dabei sind $T_1$ und $T_2$ die Rangsummen aus den jeweiligen Stichproben.
Für den Test folgt ein Prüfgröße $U$  
Dieser Test
hat keine Anforderungen über die Verteilung der Stichproben.
Somit ist keine Normalverteilung der Stichproben notwendig.


\section{Konzept Analysetool}
Das Analysetool wird mit den Logs des Fio-Tools analysieren.
Das Fio-Too durchläuft 5 - 10 mal die gleichen Tests, wo dann der stationäre Zustand, des Analysetools,
bestimmt werden soll. Das Tool wird in der Lage mit den drei Tests die eingeführt zu arbeiten.
Da die Anzahl nicht festgelegt werden kann, da die dieser Zustand an einem unbestimmten Zeitpunkt und Iteration erreicht wird,
werden verschiedene Anläufe getestet mit einer verschieden Anzahl and Fio-Tool Testanzahl.

\section{Notes}

- Analysetool mit logs aus einem HPC
- Steady State und Warm Up

----------------------------------------------------------------------

T-Test Beispiel:

Insgesamt Logs: 25

Aufgeilt in 	| Durschnitt I/O | Standardabweichung | Konfidenzintervall
Gruppe A | 13	| 43435.3	 | 3435

Gruppe B | 12	| 35254.6	 | 5456	

----------------------------------------------------------------------

Konfidenzintervall:

Beispiel es gibt insgesamt 55 logs (Das ist die Grundgesamtheit):
Stichprobe genommen sind 34

Verteilung -> normal Verteilt

Doch,

Beispiel es gibt insgesamt 55 logs (Das ist die Grundgesamtheit):
Stichprobe genommen sind 22

Verteilung -> Studentisch t Verteilt

Daraus wird der Mittelwert berechnet.

Varianz und Standardabweichung wird hier benötigt

Paper: oopsla07-georges.pdf

fsbench-tr.pdf Seite 8:

The standard deviation is a measure of how much variation there is between the
runs. The half-width of the confidence interval describes how far the true value may be
from the captured mean with a given degree of confidence (e.g., 95%). This provides a
better sense of the true mean. In addition, as more benchmark runs are performed, the
standard deviation may not decrease, but the width of the confidence interval will.
For experiments with less than 30 runs, one should be careful not to use the normal
distribution for calculating confidence intervals. This is because the central limit theorem
no longer holds with a small sample size. Instead, one must use the Student’s t-distribution.
This distribution may also be used for experiments with at least 30 runs, since in this case
it is similar to the normal distribution.

Vorstellung Programm:

- Signifikanzniveau selbst entscheidbar 0.01 bis 0.05 \\
- Wenn ein Zusammenhang zwischen den Stichproben nicht der Fall ist, sollte der p Wert aus den Proben kleiner als der HSD-Wert sein.\\
- So sollte der p Wert innerhalb des Signifikanzniveau sein und der stationäre Zustand sollte erreicht sein\\
- Die Werte sollen nicht signifikant Unterscheidbar sein, dann wäre der stationäre Zustand erreicht\\
- Wie soll ich die Logs mit den Tests abarbeiten ?\\
- Welche Logs sollen genommen werden ? 

\chapter{Verwandte Arbeiten}
\label{cha:Verwandte_Arbeiten}

Es wird in den Paper Georges et. al \cite{statistically_rigorous} mehrere Benchmarktools verwendet. 
Dies wird auch von Traeger et. al \cite{nine-year-of-bench} empfohlen um faire Tests und vergleiche machen zu können.
Festplattenbenchmarking wurde auch von Xu et. al \cite{ssd_benchmark} mit fio getesten.
Dabei wurde auf die I/O-Engine geachtet und die \textit{libaio} asynchrone I/O-Engine genutzt.
um die page cache zu bypassen und die rohe Perforamnz der Festplatten zu testen.
Dabei wird zu ein Flag für das fio gesetzt.
Der \textit{direct} Flag wurde in dieser Arbeit unverändert gelassen.
Zusätzlich wurde in Xu et. al \cite{ssd_benchmark} noch extra das Blktrace Tool (von Axboe and Brunelle 2007) genutzt. 
Für die Perforamnzanalyse wurden zwei weitere Statistiken verwendet. Die slat (submission latency) und clat (completion latency).
Für die Evaluierung wurden verschiedene Datenbanken verwendet um realitätsnahe Tests durchzuführen,
da Datenbankmanagementsysteme (DBMS) ein Beispiel für I/O-intensive Anwendungen in Rechenzentren sind.
Dabei wurden SATA HDD, SATA SSD und NVMe mit den beiden oben erwähnten Metriken der Latenz verglichen.

Das Paper Barrett et. al \cite{warmAndCold} verwendet kein Festplattenbenchmarking, aber geht auf das Benchmarking und stationären Zustand für VMs ein.
Dabei ist mit Virtuelle Maschine hier eine JIT\footnote{Just-In-Time}-Kompilierung gemeint.
Das Paper Reichelt et. al \cite{baseline_paper} verwenden ebenfalls solche VMs um diesen stationären Zustand zu ermitteln,
dabei werden ebenso statistische Tests wie der T-Test, Mann-Whitney-Test verwendet und ebenso vergleiche von Konfidenzintervallen durchgeführt. 
Wobei dort die Frage gestellt wird, wann der steady state of peak performance erreicht innerhalb einer JIT-Kompilierung.
Auch dieses Paper verwendet Methoden die in Georges et. al \cite{statistically_rigorous} empfohlen werden.

In dem Paper Chen et. al \cite{statistical_performance_pc} wird ein Beispiel Szenario vorgestellt wo die Normalverteilung nicht vorhanden ist.
Dabei werden statistische Techniken verwendet wie Normality Fitting (NNF) and Kernel Parzen Window (KPW).
Um die Normalverteilung der Leistung zu überprüfen, wird empirisch untersucht, 
ob die durch die Annahme der Normalverteilung der Ausführungszeit erhaltene Wahrscheinlichkeitsdichtefunktion
mit der tatsächlichen Wahrscheinlichkeitsdichtefunktion der Ausführungszeit übereinstimmt.
Doch in Georges et. al \cite{statistically_rigorous} im Kontext der Leistungsbewertung von Java stellt fest, dass die Leistung einzelner Benchmarks (SPECjvm98) 
auf mehreren Single-Core-Computern im Allgemeinen durch Normalverteilungen charakterisiert werden kann.
Weiterhin in Chen et. al wird der Grenzwertsatz bei solchen Messungen überprüft.
Es wird geschaut ob für welche Anzahl an Messwerten notwendig ist damit der Zentralegrenzwertsatz erfüllt.

In Alghmadi et. al \cite{when_stop_tests} wird untersucht wann ein Stop von Tests ausreicht.
Um zu wissen wann Messungen zu stoppen sind, verwendet Alghmadi et. al \cite{when_stop_tests} den mann-Whitney-rang-Test um heruaszufinden, 
wann solche Tests gestoppt werden können. Durch repetitive Muster bei der Perforamnzanalyse soll der Test dabei helfen diese zu erkennen.





\chapter{Implementierung}
\label{cha:Analysetool}
\chapter{Experimente}
\label{cha:Experimente}

\chapter{Evaluierung}
\label{cha:Evaluierung}

\chapter{Zusammenfassung und Ausblick}
\label{cha:ZuA}



% Literaturverzeichnis -----------------------------------------------------
%		Das Literaturverzeichnis wird aus der Datenbank erstellt.
%		Die genaue Verwendung von biblatex wird hier jedoch nicht erklärt.
%		Links: 	https://ctan.org/pkg/biblatex?lang=de
%						https://de.overleaf.com/learn/latex/Articles/Getting_started_with_BibLaTeX
% --------------------------------------------------------------------------

\printbibliography

% \setcounter{page}{122}
% \pagenumbering{gobble}
%\pagenumbering{gobble}
\addchap{Erklärung}
Ich versichere, dass ich die vorliegende Arbeit mit dem Thema:

\begin{center}
\textit{\glqq\titel\grqq}\\[1em]
\end{center}
			
selbständig und nur unter Verwendung der angegebenen Quellen und Hilfsmittel angefertigt habe, insbesondere sind wörtliche oder sinngemäße Zitate als solche gekennzeichnet. Mir ist bekannt, dass Zuwiderhandlung auch nachträglich zur Aberkennung des Abschlusses führen kann. Ich versichere, dass das elektronische Exemplar mit den gedruckten Exemplaren übereinstimmt.
\par
\ort, den \eingereicht


\rule[-0.2cm]{5cm}{0.5pt}

\textsc{\autor} 
	% Selbständigkeitserklärung 

% Anhang -------------------------------------------------------------------
%		Die Inhalte des Anhangs werden analog zu den Kapiteln inkludiert.
%		Dies geschieht in der Datei Anhang.tex
% --------------------------------------------------------------------------
\appendix
\clearpage
\renewcommand*{\thesection}{\Alph{section}} 
\pagenumbering{Roman}
%\include{Inhalt/Anhang}



% Index --------------------------------------------------------------------
%		Zum Erstellen eines Index, die folgende Zeile auskommentieren.
% --------------------------------------------------------------------------
%\printindex		% Index hier einfügen
%\ofoot{}
%\include{Inhalt/Thesen}	% Thesen 

\end{document}
